\chapter{Statistical Estimation}
\chaptermark{Statistical Estimation}
\label{chap:estimation}

\begingroup
\clubpenalty=10000
\widowpenalty=10000
\newenvironment{estimationbox}[1]{\begin{tcolorbox}[
  colback=CourseBlue!3!white,colframe=CourseBlue!65!white,
  colbacktitle=CourseBlue!12!white,coltitle=CourseDark,
  fonttitle=\sffamily\bfseries,title={#1},boxrule=0.5pt,
  arc=0pt,left=7pt,right=7pt,top=5pt,bottom=5pt]}{\end{tcolorbox}}
% Shared analytic/simulation data charts for the conspect and the Beamer deck.
% Callers may override dimensions and the data path before loading this file.
\providecommand{\EstDataPath}{figures/data}
\providecommand{\EstPlotWidth}{11.8cm}
\providecommand{\EstPlotHeight}{4.8cm}
\definecolor{EstBlue}{HTML}{1F4E79}
\definecolor{EstOrange}{HTML}{C46A13}
\definecolor{EstTeal}{HTML}{007D80}
\definecolor{EstMuted}{HTML}{526477}
\pgfplotsset{est axis/.style={
  scale only axis,width=\EstPlotWidth,height=\EstPlotHeight,
  axis lines=left,enlargelimits=false,
  grid=major,grid style={black!9},
  tick label style={font=\footnotesize},label style={font=\footnotesize},
  legend style={draw=none,fill=none,font=\footnotesize,
    at={(0.5,1.05)},anchor=south,legend columns=-1,column sep=0.8em},
  every axis plot/.append style={no marks,line width=1.15pt},
}}

\newcommand{\EstLikelihoodPlot}{%
\begin{tikzpicture}\begin{axis}[est axis,
  xmin=0,xmax=1.2,ymin=0,ymax=1.06,
  xlabel={Candidate rate $\lambda$ (per minute)},ylabel={Relative likelihood},
  xtick={0,0.2,0.4,0.6,0.8,1,1.2}]
  \addplot[EstBlue] table[x=rate,y=n5] {\EstDataPath/ch5-likelihood.dat};
  \addlegendentry{$n=5$, total wait $15$}
  \addplot[EstOrange,dashed] table[x=rate,y=n50] {\EstDataPath/ch5-likelihood.dat};
  \addlegendentry{$n=50$, total wait $150$}
  \addplot[EstMuted,densely dotted,forget plot] coordinates {(0.333333,0) (0.333333,1)};
\end{axis}\end{tikzpicture}%
}

\newcommand{\EstUniformPlot}{%
\begin{tikzpicture}\begin{axis}[est axis,
  xmin=0,xmax=2,ymin=0,ymax=5.4,xlabel={Estimated upper endpoint},ylabel={Density},
  xtick={0,0.5,1,1.5,2}]
  \addplot[EstBlue,const plot] table[x=endpoint,y=moments] {\EstDataPath/ch5-uniform-estimates.dat};
  \addlegendentry{Moment estimator $2\bar X$}
  \addplot[EstOrange,const plot] table[x=endpoint,y=mle] {\EstDataPath/ch5-uniform-estimates.dat};
  \addlegendentry{MLE $\max_iX_i$}
  \addplot[EstMuted,densely dotted,forget plot] coordinates {(1,0) (1,5.2)};
  \node[anchor=north west,font=\footnotesize,text=EstMuted,fill=white,inner sep=2pt]
    at (axis cs:1.06,4.95) {True endpoint: $1$};
\end{axis}\end{tikzpicture}%
}

\newcommand{\EstVarianceLossPlot}{%
\begin{tikzpicture}\begin{axis}[est axis,
  xmin=0.5,xmax=6,ymin=-1,ymax=8.2,
  xlabel={Predicted standard deviation $\sigma$},ylabel={Loss contribution},
  xtick={1,2,3,4,5,6}]
  \addplot[EstBlue] table[x=sigma,y=total] {\EstDataPath/ch5-variance-loss.dat};
  \addlegendentry{Total: $\log\sigma+2/\sigma^2$}
  \addplot[EstOrange,dashed] table[x=sigma,y=scaled_residual] {\EstDataPath/ch5-variance-loss.dat};
  \addlegendentry{$2/\sigma^2$}
  \addplot[EstTeal,densely dotted] table[x=sigma,y=log_scale] {\EstDataPath/ch5-variance-loss.dat};
  \addlegendentry{$\log\sigma$}
  \addplot[EstBlue,only marks,mark=*,mark size=2pt,forget plot] coordinates {(2,1.19314718)};
  \node[font=\footnotesize,anchor=north,text=EstBlue,fill=white,inner sep=2pt]
    at (axis cs:2,0.8) {Minimum at $\sigma=2$};
\end{axis}\end{tikzpicture}%
}

\newcommand{\EstBetaPlot}{%
\begin{tikzpicture}\begin{axis}[est axis,
  xmin=0,xmax=1,ymin=0,ymax=3.4,
  xlabel={Bernoulli probability $p$},ylabel={Density over $p$},
  xtick={0,0.2,0.4,0.6,0.8,1}]
  \addplot[EstMuted,dashed] table[x=p,y=prior] {\EstDataPath/ch5-beta.dat};
  \addlegendentry{Prior: Beta$(2,2)$}
  \addplot[EstBlue] table[x=p,y=posterior] {\EstDataPath/ch5-beta.dat};
  \addlegendentry{Posterior: Beta$(9,5)$}
  \addplot[EstOrange,densely dotted,forget plot] coordinates {(0.6666667,0) (0.6666667,3.15)};
\end{axis}\end{tikzpicture}%
}

\newcommand{\EstGaussianPriorPlot}{%
\begin{tikzpicture}\begin{axis}[est axis,
  xmin=-3,xmax=5,ymin=0,ymax=0.61,
  xlabel={Candidate mean $\mu$},ylabel={Density / rescaled likelihood},
  xtick={-2,0,0.8,1.6,3,5},xticklabels={$-2$,$0$,$0.8$,$1.6$,$3$,$5$}]
  \addplot[EstMuted,dashed] table[x=mu,y=prior] {\EstDataPath/ch5-gaussian-prior.dat};
  \addlegendentry{Prior}
  \addplot[EstOrange,dashdotted] table[x=mu,y=likelihood] {\EstDataPath/ch5-gaussian-prior.dat};
  \addlegendentry{Rescaled likelihood}
  \addplot[EstBlue] table[x=mu,y=posterior] {\EstDataPath/ch5-gaussian-prior.dat};
  \addlegendentry{Posterior}
\end{axis}\end{tikzpicture}%
}


\section{An unknown parameter in a probability model}
\label{sec:estimation-model}

Suppose five independent waits between requests are
\[
  t_1,\ldots,t_5=1,2,3,4,5\ \text{minutes}.
\]
Their average is three minutes. In Chapter~\ref{chap:limit-theorems},
we specified an exponential rate and studied how sample averages fluctuate.
Now we keep the exponential model but treat its rate \(\lambda>0\) as
unknown. Which rate should we report from these observations?

A \emph{statistical model} is a family of possible observation laws,
written \(\{P_\theta:\theta\in\Theta\}\). The parameter \(\theta\)
can be a scalar or a vector, and \(\Theta\) specifies its allowed values.
For exponential waits, \(\theta=\lambda\) and \(\Theta=(0,\infty)\).
For a Gaussian, the unknown parameter might be \(\theta=(\mu,\sigma^2)\),
with \(\mu\in\Real\) and \(\sigma^2>0\). Choosing a family is an
assumption about the data-generating process; estimating its parameters
does not establish that the family is correct.

\begin{estimationbox}{Definition: estimator and estimate}
An \emph{estimator} is a rule \(\widehat\theta=T(X_1,\ldots,X_n)\)
that uses the observations to estimate a parameter. It is random before
we collect the data. Its realized value \(T(x_1,\ldots,x_n)\) is an
\emph{estimate}. The rule cannot use the unknown true parameter.
\end{estimationbox}

We retain the sample notation from Section~\ref{sec:samples}. For the waits,
\(1/\bar T_n\) is a possible estimator, whereas \(1/3\) requests per minute
is its value on this dataset. A point estimate supplies a single parameter
value. Quantifying its uncertainty is a further question.

The observations must also distinguish the parameter values of interest.
A model is \emph{identifiable} when \(P_\theta=P_{\theta'}\) implies
\(\theta=\theta'\). For example, if \(X\sim\mathcal N(a^2,1)\) and
\(a\in\Real\), the values \(a\) and \(-a\) give exactly the same
observation law. No amount of such data identifies the sign. Restricting
\(a\ge0\), or estimating \(a^2\), removes this particular ambiguity.

\section{Method of moments}
\label{sec:method-moments}

The model links its parameters to population summaries that we already know
how to estimate by averaging. The method of moments solves those links in
reverse. For exponential waits, the known relation
\(\Expect_\lambda[T]=1/\lambda\) says that a large rate corresponds to a
short average wait. Replacing the population mean by the observed mean gives
\[
  \frac1{\widehat\lambda_{\mathrm{MM}}}=\bar t_n,
  \qquad \widehat\lambda_{\mathrm{MM}}=\frac1{\bar t_n}.
\]
The five waits therefore give \(\widehat\lambda_{\mathrm{MM}}=1/3\)
requests per minute. The units also follow from inverting a mean measured
in minutes.

\begin{estimationbox}{Definition: method of moments}
For a model with \(k\) unknown parameters, choose \(k\) population moments
\(m_j(\theta)=\Expect_\theta[X^j]\) and their empirical counterparts
\(\widehat m_j=n^{-1}\sum_i X_i^j\). A method-of-moments estimator solves
\[
  m_j(\widehat\theta_{\mathrm{MM}})=\widehat m_j,
  \qquad j=1,\ldots,k,
\]
with a solution in \(\Theta\). Other chosen functions of \(X\) can replace
the powers when their expectations identify the parameter.
\end{estimationbox}

For Bernoulli data the mean is \(p\), so matching means gives
\(\widehat p_{\mathrm{MM}}=\bar X_n\). For Poisson data the mean is
\(\lambda\), so \(\widehat\lambda_{\mathrm{MM}}=\bar X_n\).
These are applications of the distribution moments from
Chapter~\ref{chap:distributions}, not new assumptions about those moments.

\subsection{Two moments for a Gamma model}

Suppose instead that independent positive observations follow a Gamma
distribution with shape \(\alpha>0\) and \emph{rate} \(\beta>0\).
The earlier formulas are
\[
  \Expect[X]=\frac\alpha\beta,\qquad
  \Var(X)=\frac\alpha{\beta^2}.
\]
Let \(\bar x=n^{-1}\sum_i x_i\) and
\(v_n=n^{-1}\sum_i(x_i-\bar x)^2=\widehat m_2-\widehat m_1^2\).
Matching the first two raw moments is equivalent to solving
\[
  \bar x=\frac{\widehat\alpha}{\widehat\beta},\qquad
  v_n=\frac{\widehat\alpha}{\widehat\beta^2}.
\]
Dividing the second equation by the first gives
\begin{equation}
  \widehat\beta_{\mathrm{MM}}=\frac{\bar x}{v_n},\qquad
  \widehat\alpha_{\mathrm{MM}}=\frac{\bar x^2}{v_n}.
  \label{eq:gamma-moments-estimate}
\end{equation}
For the values \(1,2,3,4,5\), \(\bar x=3\) and \(v_n=2\), giving
\(\widehat\alpha=4.5\) and \(\widehat\beta=1.5\) per minute.
The Gamma fit can match both the mean and spread. An exponential fit has
only one free parameter and ties the variance to the square of the mean.

The denominator here is \(n\) because we are matching empirical raw
moments. Replacing \(v_n\) by the variance with denominator \(n-1\)
would give a different rule. If every observation is identical, \(v_n=0\)
and these equations give no finite positive Gamma parameters.

\subsection{When matching moments is justified}

The required population moments must exist, and the chosen moments must
distinguish the parameters. Even then, a finite dataset can give multiple
solutions, no admissible solution, or an unstable solution. For example,
dividing by a very small empirical variance can produce very large estimates.

The LLN from Section~\ref{sec:lln} explains why the method often works:
empirical moments approach their population values. If the inverse map from
moments to parameters is well-defined and continuous near the true moments,
the estimates approach the true parameters. This argument needs those
conditions; solving equations alone does not guarantee accurate estimation.

\section{Likelihood and maximum likelihood}
\label{sec:likelihood-estimation}

Matching moments uses selected summaries. Likelihood uses the joint
probability model of the observations. We now hold the recorded data fixed
and vary the candidate parameter.

\begin{estimationbox}{Definition: likelihood and maximum likelihood estimator}
For i.i.d. observations with PMF or density \(f_\theta\), the likelihood is
\[
  L(\theta;x_1,\ldots,x_n)=\prod_{i=1}^n f_\theta(x_i).
\]
A maximum likelihood estimator (MLE) is a parameter value attaining the
largest likelihood within the allowed parameter space:
\[
  \widehat\theta_{\mathrm{ML}}\in\argmax_{\theta\in\Theta}L(\theta;x).
\]
The maximum need not be unique, and it need not exist.
\end{estimationbox}

For discrete observations, the likelihood is the probability of the recorded
sample, viewed as a function of \(\theta\). For continuous observations it
is a joint \emph{density} value. Comparing densities compares probabilities
of small neighborhoods of the recorded observations to first order, where
the neighborhood sizes do not depend on \(\theta\). Likelihood is not a
probability distribution over parameter values, and it need not integrate
to one as a function of \(\theta\).

If observations are dependent, we must use their actual joint law. Multiplying
marginal densities in that case generally gives the wrong likelihood.

\subsection{Log-likelihood and the exponential rate}

The logarithm is strictly increasing, so maximizing a positive likelihood
is equivalent to maximizing its log. Under independence,
\[
  \ell(\theta;x)=\log L(\theta;x)=\sum_{i=1}^n\log f_\theta(x_i).
\]
Adding logarithms is also numerically safer than multiplying many small
densities. A factor independent of \(\theta\) can be dropped when locating
the maximum; a factor that depends on \(\theta\) cannot.

For positive exponential waits, let \(s=\sum_i t_i>0\). Then
\[
  L(\lambda)=\lambda^n e^{-\lambda s},\qquad
  \ell(\lambda)=n\log\lambda-\lambda s.
\]
Differentiation gives
\[
  \ell'(\lambda)=\frac n\lambda-s,\qquad
  \ell''(\lambda)=-\frac n{\lambda^2}<0.
\]
Thus the unique maximum is
\begin{equation}
  \widehat\lambda_{\mathrm{ML}}=\frac n{\sum_i t_i}=\frac1{\bar t_n}.
  \label{eq:exponential-mle}
\end{equation}
Here maximum likelihood and method of moments agree. For our five waits,
both report \(1/3\) per minute. Agreement in this model does not mean the
two methods are equivalent in general.

\begin{figure}[H]
  \centering\EstLikelihoodPlot
  \caption{Relative likelihood \(L(\lambda)/L(\widehat\lambda)\) for two
  illustrative datasets with the same average wait of three minutes:
  \(n=5,s=15\), and \(n=50,s=150\). Both maximize at \(1/3\), but the
  larger sample distinguishes nearby rates more sharply. These curves are
  likelihood ratios, not densities over the rate.}
  \label{fig:est-likelihood}
\end{figure}

\subsection{The support can determine the maximum}

Consider a different model, \(X_i\sim\operatorname{Uniform}(0,\theta)\)
with \(\theta>0\), using density
\(f_\theta(x)=\theta^{-1}\mathbf1_{\{0\le x\le\theta\}}\).
For a positive sample maximum \(x_{(n)}=\max_i x_i\), the joint likelihood is
\[
  L(\theta)=\theta^{-n}\mathbf1_{\{\theta\ge x_{(n)}\}}.
\]
Any upper endpoint below an observed value has likelihood zero. On the
admissible region, the likelihood decreases with \(\theta\), so
\[
  \widehat\theta_{\mathrm{ML}}=x_{(n)},\qquad
  \widehat\theta_{\mathrm{MM}}=2\bar x_n,
\]
where the moment estimate follows from \(\Expect[X]=\theta/2\).
For \(x=(0.2,0.4,0.5,0.7,0.9)\), the two estimates are \(0.9\) and
\(1.08\). The maximum occurs at a support boundary. Solving
\(\ell'(\theta)=0\) would miss it entirely.

More generally, a stationary point need not be a global maximum. Parameter
constraints, boundaries, and unbounded likelihoods must be checked. For
instance, fitting a Gaussian with free variance to identical observations
lets the likelihood increase without bound as the variance approaches zero.

\section{Fitting a Gaussian mean and variance}
\label{sec:gaussian-estimation}

Let \(X_i\) be i.i.d. \(\mathcal N(\mu,v)\), writing \(v=\sigma^2>0\)
to keep variance distinct from standard deviation. The log-likelihood is
\[
  \ell(\mu,v)=-\frac n2\log(2\pi)-\frac n2\log v
              -\frac1{2v}\sum_{i=1}^n(x_i-\mu)^2.
\]
For each fixed \(v\), maximizing over \(\mu\) minimizes the squared
deviations. The identity
\[
  \sum_i(x_i-\mu)^2=\sum_i(x_i-\bar x)^2+n(\mu-\bar x)^2
\]
shows directly that \(\widehat\mu_{\mathrm{ML}}=\bar x\). With
\(Q=\sum_i(x_i-\bar x)^2>0\), differentiate with respect to \(v\):
\[
  \frac{\partial\ell(\bar x,v)}{\partial v}
  =-\frac n{2v}+\frac Q{2v^2}
  =\frac{Q-nv}{2v^2}.
\]
This derivative is positive for \(v<Q/n\) and negative for \(v>Q/n\).

\begin{estimationbox}{Result: Gaussian maximum likelihood estimates}
For a nonconstant i.i.d. Gaussian sample,
\[
  \widehat\mu_{\mathrm{ML}}=\bar X_n,\qquad
  \widehat v_{\mathrm{ML}}=\frac1n\sum_i(X_i-\bar X_n)^2.
\]
If all observations coincide, there is no finite maximizer with \(v>0\).
\end{estimationbox}

The method of moments gives these same formulas by matching
\(\Expect[X]=\mu\) and \(\Expect[X^2]=\mu^2+v\). The two methods agree
again, for a different reason than a general equivalence. The mean and
variance formulas still make sense for non-Gaussian data, but their Gaussian
maximum-likelihood interpretation then depends on using a Gaussian model.

\section{Bias, variance, and consistency}
\label{sec:estimator-quality}

To evaluate a rule, imagine repeating the entire sampling experiment with
the true parameter \(\theta\) fixed. The subscript in \(\Expect_\theta\)
or \(\Var_\theta\) specifies that sampling distribution.

\begin{estimationbox}{Definition: bias and mean squared error}
For a scalar estimator with finite second moment,
\[
  \operatorname{Bias}_\theta(\widehat\theta)=\Expect_\theta[\widehat\theta]-\theta,
  \qquad
  \operatorname{MSE}_\theta(\widehat\theta)
  =\Expect_\theta[(\widehat\theta-\theta)^2].
\]
An estimator is \emph{unbiased} if its bias is zero for every allowed
parameter value. Expanding about \(\Expect_\theta[\widehat\theta]\) gives
\[
  \operatorname{MSE}_\theta(\widehat\theta)
  =\Var_\theta(\widehat\theta)+\operatorname{Bias}_\theta(\widehat\theta)^2.
\]
The cross term vanishes because the centered estimator has expectation zero.
\end{estimationbox}

Bias is an average error across repeated datasets; it is not the error in
one realized estimate. An unbiased estimator can fluctuate substantially.
Under squared-error comparison, reducing variance can compensate for some bias.

\subsection{Why the sample variance uses \(n-1\)}

This answers the question left open in Section~\ref{sec:samples}.
For i.i.d. observations with mean \(\mu\) and finite variance \(\sigma^2\),
\[
  \sum_i(X_i-\bar X_n)^2
  =\sum_i(X_i-\mu)^2-n(\bar X_n-\mu)^2.
\]
Taking expectations and using \(\Var(\bar X_n)=\sigma^2/n\) gives
\[
  \Expect\!\left[\sum_i(X_i-\bar X_n)^2\right]
  =n\sigma^2-n\frac{\sigma^2}{n}=(n-1)\sigma^2.
\]
Therefore
\begin{equation}
  \Expect[\widehat v_{\mathrm{ML}}]=\frac{n-1}{n}\sigma^2,
  \qquad
  \Expect[S_n^2]=\sigma^2,\quad
  S_n^2=\frac1{n-1}\sum_i(X_i-\bar X_n)^2.
  \label{eq:variance-bias-correction}
\end{equation}
Centering at the fitted mean removes part of the observed variation.
The \(n-1\) correction makes the variance estimator unbiased for \(n\ge2\).
This expectation calculation needs no Gaussian assumption. Under a Gaussian
model the denominator \(n\) still maximizes the likelihood; unbiasedness
and maximum likelihood are different criteria.

\subsection{A biased estimate can have smaller squared error}

For independent \(\operatorname{Uniform}(0,\theta)\) observations,
the moment estimator \(2\bar X_n\) is unbiased, with MSE
\(\theta^2/(3n)\). For \(M=\max_iX_i\),
\(\Prob_\theta(M\le m)=(m/\theta)^n\) on \([0,\theta]\).
Its density is \(nm^{n-1}/\theta^n\); integrating \(m\) and \(m^2\)
against it gives
\[
  \Expect_\theta[M]=\frac n{n+1}\theta,\qquad
  \Expect_\theta[M^2]=\frac n{n+2}\theta^2.
\]
Consequently its bias is \(-\theta/(n+1)\) and
\[
  \operatorname{MSE}_\theta(M)
  =\Expect_\theta[M^2]-2\theta\Expect_\theta[M]+\theta^2
  =\frac{2\theta^2}{(n+1)(n+2)}.
\]
At \(n=5,\theta=1\), these MSEs are \(1/15\approx0.0667\) for
method of moments and \(1/21\approx0.0476\) for the biased MLE.
This comparison concerns these two rules under this model and squared error;
it does not make the MLE universally best.

\begin{figure}[H]
  \centering\EstUniformPlot
  \caption{Estimates from 50,000 independent datasets, each containing five
  \(\operatorname{Uniform}(0,1)\) observations. The moment estimator spreads
  on both sides of the true endpoint. The maximum always underestimates it,
  but its smaller spread yields the smaller MSE in this example. Curves
  show density-normalized histogram steps on common bins.}
  \label{fig:est-uniform}
\end{figure}

\subsection{Consistency concerns growing sample size}

\begin{estimationbox}{Definition: consistency}
A sequence of estimators is \emph{consistent} for \(\theta\) if for every
\(\varepsilon>0\),
\[
  \Prob_\theta(|\widehat\theta_n-\theta|>\varepsilon)\longrightarrow0
  \quad\text{as }n\to\infty.
\]
This is convergence in probability of the estimation error to zero.
\end{estimationbox}

For exponential waits, the LLN gives \(\bar T_n\to1/\lambda\) in
probability. Taking reciprocals is continuous near this positive limit,
so \(\widehat\lambda=n/\sum_iT_i\to\lambda\). This proves consistency
without implying unbiasedness. In fact, for \(n>1\), the Gamma law of
\(\sum_iT_i\) gives
\[
  \Expect_\lambda[\widehat\lambda]=\frac n{n-1}\lambda.
\]
The rate estimate has positive finite-sample bias, which tends to zero.
For \(n=1\) its expectation is infinite, although the estimate itself is
finite almost surely. Conversely, using just \(X_1\) as an estimate of
the population mean is unbiased but generally inconsistent: its spread
does not decrease when more observations become available.

\section{Likelihood gives a training loss}
\label{sec:likelihood-loss}

So far one distribution described every observation. For prediction, specify
a distribution of the response \(Y_i\) conditional on its observed features
\(x_i\). Treat the feature values as fixed and assume responses are
conditionally independent. The conditional likelihood is
\(\prod_i p_\theta(y_i\mid x_i)\). Estimation chooses \(\theta\) to make
this observed response sample have large conditional likelihood.

\subsection{Gaussian regression and squared error}

Let \(Y_i\mid x_i\sim\mathcal N(f_\theta(x_i),\sigma^2)\), with the
same known \(\sigma^2>0\) for every response. Then
\begin{equation}
  -\ell(\theta)=\frac n2\log(2\pi\sigma^2)
     +\frac1{2\sigma^2}\sum_i\bigl(y_i-f_\theta(x_i)\bigr)^2.
  \label{eq:gaussian-regression-loss}
\end{equation}
The first term and the positive scale factor do not change which
\(\theta\) minimizes the expression. Thus maximizing this likelihood is
equivalent to minimizing mean squared error. This explains a familiar loss
through an explicit noise model. Chapter~\ref{chap:joint-distributions}
showed why a conditional mean predicts under squared error; here we estimate
the parameters of a chosen conditional-mean model from data.

If a common variance \(v\) is also unknown, keep the term
\(\tfrac n2\log v\). For any fixed \(\theta\) with positive residual
sum of squares, its fitted value is
\(\widehat v(\theta)=n^{-1}\sum_i(y_i-f_\theta(x_i))^2\).
When the model fits every response exactly, this joint likelihood is
unbounded as \(v\downarrow0\). A finite solution then requires an
additional constraint or a different model.

\subsection{A predicted variance changes the objective}

Suppose the model predicts both a mean \(\mu_\theta(x_i)\) and a positive
variance \(v_\theta(x_i)\). Up to the constant \(\tfrac n2\log(2\pi)\),
\begin{equation}
  -\ell(\theta)=\frac12\sum_i\left[
    \log v_\theta(x_i)+\frac{(y_i-\mu_\theta(x_i))^2}{v_\theta(x_i)}
  \right].
  \label{eq:heteroscedastic-loss}
\end{equation}
The inverse variance weights the squared residual, but predicting a large
variance also increases the logarithmic term. Dropping that term would let
the model hide every residual by making its variance arbitrarily large.
A parameterization such as \(v_\theta(x)=e^{s_\theta(x)}\) enforces
positivity. It does not by itself prevent variance collapse at a perfectly
fitted observation.

\begin{figure}[H]
  \centering\EstVarianceLossPlot
  \caption{For one fixed residual \(y-\mu=2\), Gaussian negative
  log-likelihood, excluding its constant, is \(\log\sigma+2/\sigma^2\).
  Its minimum occurs at \(\sigma=2\). The logarithmic term prevents
  the scale from growing without cost. For residual zero, in contrast,
  the loss tends to \(-\infty\) as \(\sigma\downarrow0\).}
  \label{fig:est-variance-loss}
\end{figure}

\subsection{Bernoulli responses and binary log loss}

For \(Y_i\in\{0,1\}\), let the model predict
\(p_i=p_\theta(x_i)\in(0,1)\). The conditional PMF is
\(p_i^{y_i}(1-p_i)^{1-y_i}\), so maximizing likelihood minimizes
\begin{equation}
  -\ell(\theta)
  =-\sum_i\left[y_i\log p_i+(1-y_i)\log(1-p_i)\right].
  \label{eq:bernoulli-log-loss}
\end{equation}
For an observed positive response, assigning probability \(0.9\) contributes
\(-\log0.9\approx0.105\), whereas probability \(0.1\) contributes
\(-\log0.1\approx2.303\). Confident wrong predictions incur a large cost.
The information-theoretic interpretation of this loss belongs to the
later chapter on entropy and cross-entropy.

The statistical model determines the objective. An optimization algorithm
searches for parameters that minimize it; a numerical stationary point need
not be a global optimum. Good training likelihood also does not guarantee
good prediction on new data. Evaluation requires the independent-data
reasoning from Section~\ref{sec:monte-carlo}.

\section{MAP estimation and prior information}
\label{sec:map-estimation}

Maximum likelihood compares how well candidate parameters account for the
observations. A Bayesian model also specifies a prior distribution on the
parameter before using the current data. A prior density \(\pi(\theta)\)
describes that parameter uncertainty. Conditional on a parameter value,
the same observation model supplies the likelihood.

The conditioning rules from Chapters~\ref{chap:probability-models} and
\ref{chap:joint-distributions} give
\[
  \pi(\theta\mid x)
  =\frac{L(\theta;x)\pi(\theta)}
         {\int_\Theta L(u;x)\pi(u)\,du},
\]
provided the denominator is positive and finite. For discrete parameters,
replace the integral by a sum. The denominator is constant in \(\theta\),
so it does not affect the location of a posterior maximum.

\begin{estimationbox}{Definition: maximum a posteriori estimate}
A MAP estimate maximizes the posterior density (or posterior mass for
a discrete parameter):
\[
  \widehat\theta_{\mathrm{MAP}}\in
  \argmax_{\theta\in\Theta}\{\ell(\theta;x)+\log\pi(\theta)\}.
\]
Equivalently, it minimizes negative log-likelihood plus negative log-prior.
This is a point summary of the posterior, not the entire posterior distribution.
\end{estimationbox}

\subsection{A Bernoulli probability with a Beta prior}

Take \(k\) successes in \(n\) independent Bernoulli observations.
The likelihood is proportional to \(p^k(1-p)^{n-k}\); its MLE is
\(k/n\), including boundary values if \(p\in[0,1]\).
To express prior information about a probability, introduce the Beta family:
\[
  \pi(p)=\frac{p^{a-1}(1-p)^{b-1}}{B(a,b)},\quad 0<p<1,\quad a,b>0,
  \qquad B(a,b)=\int_0^1u^{a-1}(1-u)^{b-1}\,du.
\]
Here \(a,b\) are specified prior parameters. Multiplying prior and likelihood
adds their powers, yielding
\[
  p\mid x\sim\operatorname{Beta}(a+k,b+n-k).
\]
For \(a+k>1\) and \(b+n-k>1\), differentiation of the log posterior gives
\begin{equation}
  \widehat p_{\mathrm{MAP}}=\frac{k+a-1}{n+a+b-2}.
  \label{eq:beta-map}
\end{equation}
The conditions ensure an interior mode. Outside them, check the boundaries
instead of applying this formula blindly.

With a \(\operatorname{Beta}(2,2)\) prior and seven successes in ten trials,
the posterior is \(\operatorname{Beta}(9,5)\), and
\[
  \widehat p_{\mathrm{ML}}=\frac7{10}=0.7,\qquad
  \widehat p_{\mathrm{MAP}}=\frac8{12}=\frac23.
\]
The symmetric prior favors probabilities closer to \(1/2\), moving this
estimate slightly downward. A different prior can move it differently.

\begin{figure}[H]
  \centering\EstBetaPlot
  \caption{The \(\operatorname{Beta}(2,2)\) prior and
  \(\operatorname{Beta}(9,5)\) posterior after seven successes in ten trials.
  Both are normalized densities over \(p\). The posterior mode is \(2/3\);
  it summarizes the distribution without describing its whole spread.}
  \label{fig:est-beta}
\end{figure}

The posterior mean is another point summary. For a Beta\((A,B)\) law,
inserting one extra factor of \(p\) into the normalizing integral gives
\[
  \Expect[p]=\frac{B(A+1,B)}{B(A,B)}=\frac{A}{A+B}.
\]
Here it is
\(9/14\approx0.643\), different from the MAP. The conditional-mean result
in Section~\ref{sec:conditional-moments} makes the posterior mean the
minimizer of posterior expected squared error. The mode answers a different
question. For continuous parameters, MAP also depends on the chosen
parameterization because a transformed density includes a Jacobian factor.

\subsection{A Gaussian prior pulls an estimated mean toward its center}

Let \(X_i\sim\mathcal N(\mu,\sigma^2)\) with known \(\sigma^2>0\),
and use prior \(\mu\sim\mathcal N(m_0,\tau^2)\), \(\tau^2>0\).
Here \(\sigma^2\) describes variation in observations, whereas \(\tau^2\)
describes prior uncertainty about their mean. Terms depending on \(\mu\)
in negative log posterior are
\[
  \frac1{2\sigma^2}\sum_i(x_i-\mu)^2+
  \frac1{2\tau^2}(\mu-m_0)^2.
\]
Setting its derivative to zero gives
\begin{equation}
  \widehat\mu_{\mathrm{MAP}}
  =\frac{(n/\sigma^2)\bar x+(1/\tau^2)m_0}{n/\sigma^2+1/\tau^2}.
  \label{eq:gaussian-map}
\end{equation}
Completing the square shows that the posterior is Gaussian with this mean
and variance \((n/\sigma^2+1/\tau^2)^{-1}\). Its mean and mode coincide
in this special case. The weights are precisions, meaning inverse variances:
the variance of \(\bar X_n\) is \(\sigma^2/n\), and the prior variance
is \(\tau^2\). As \(n\) grows with the prior fixed, the data receive
more relative weight.

For \(n=4,\sigma^2=4,\bar x=1.6,m_0=0,\tau^2=1\), both precisions
are one. The MLE is \(1.6\), the MAP is \(0.8\), and the posterior
variance is \(1/2\).

\begin{figure}[H]
  \centering\EstGaussianPriorPlot
  \caption{Prior \(\mathcal N(0,1)\), likelihood as a function of \(\mu\)
  rescaled to integrate to one, and posterior \(\mathcal N(0.8,0.5)\) in
  the stated example. The rescaled likelihood has shape \(\mathcal N(1.6,1)\);
  its normalization is only for the plot and does not turn it into a prior
  or posterior.}
  \label{fig:est-gaussian-prior}
\end{figure}

\subsection{Gaussian priors and quadratic regularization}

In Gaussian regression with known observation variance \(\sigma^2\),
put a prior \(w\sim\mathcal N(0,\tau^2I)\) on the model coefficients.
MAP minimizes, up to constants,
\[
  \frac1{2\sigma^2}\sum_i(y_i-f_w(x_i))^2+
  \frac1{2\tau^2}\|w\|^2.
\]
Multiplying by \(2\sigma^2\) gives the equivalent objective
\begin{equation}
  \sum_i(y_i-f_w(x_i))^2+\alpha\|w\|^2,
  \qquad \alpha=\frac{\sigma^2}{\tau^2}.
  \label{eq:map-quadratic-penalty}
\end{equation}
For a linear model this is ridge regression, whose linear-algebra solution
comes later in the course. If the data term is an \emph{average} squared
error, the same prior gives coefficient \(\alpha/n\), not \(\alpha\).
The normalization of the loss matters when interpreting a penalty.
An unpenalized intercept would require a different prior choice for that
coefficient. More generally, MAP imposes the penalty \(-\log\pi(w)\);
a quadratic penalty specifically corresponds to a Gaussian prior.

Method of moments matches selected population summaries to sample summaries.
MLE maximizes the observation likelihood. MAP adds an explicit prior on the
parameter. Their estimates can coincide, but their definitions and assumptions
remain different. None of these point-estimation rules by itself supplies
a confidence interval or a guarantee of predictive accuracy.
\endgroup
